\section{Low-rank moments and mixing forms}\label{sec:moments}

We prepare point moments for use as witnesses in the hole equations.
The gradient identities must hold on the whole cut space.  The
low-rank decomposition below will be used to reduce the obstructions
to those identities to prescribed scalar tests.  A separate bound
on sparse label supports of projected pin directions identifies a
bounded set of labels to avoid, so that the witness directions remain
independent of the pins.  We then choose the mixing matrices to control evaluations
at one point and at pairs of points with distinct labels.

The pin budget $K$ is fixed before the parameters chosen in this
section.  Set
\begin{equation}\label{eq:moment-parameters}
 R_*=2(K+20),\qquad
 j_*\ge 10(K+1)(R_*+20).
\end{equation}
All inequalities concerning ranks in this section are over $\F$.
The selector size $b$ will be chosen below.  Put
$p_*=\sum_{j=0}^{j_*}\binom bj$ and write $d_{\mathrm b}=(g^2+3)n$ for the
number of ordinary base coordinates.  Thus
$\cB=\F^{p_*}\otimes\F^{1+d_{\mathrm b}}$ and
$v(s,z)=p_s\otimes(1,z)$, as in \Cref{sec:geometry}.
The constant base coordinate has index $0$.

\subsection{Boolean moment matrices}

The low-rank decomposition below uses the quotient and multiplication
operators familiar from flat extensions of moment matrices; see
Laurent and Mourrain~\citep[Section~2, especially Lemma~2.1]{laurentmourrain2009}.
Here we work over $\F$ with separate selector and base degree bounds.
The quotient pairing is nondegenerate, and no positivity is used.

\begin{lemma}\label{lem:base-moments}
The span of $(1,z)(1,z)\trans$, $z\in\F^{d_{\mathrm b}}$, consists precisely of
symmetric $(1+d_{\mathrm b})$-by-$(1+d_{\mathrm b})$ matrices $Z$ satisfying
$Z_{ii}=Z_{0i}$ for every $i$.
For any $u$ distinct selectors, their evaluation vectors $p_s$ are
linearly independent if $j_*\ge u-1$.
\end{lemma}
\begin{proof}
Every point matrix has the stated identities.  Conversely the space
specified by these identities has a basis consisting of the matrix
supported at $(0,0)$, the matrices supported at
$(0,i),(i,0),(i,i)$, and the matrices supported at $(i,j),(j,i)$ for
$0<i<j$.  The first is the point matrix at zero.  The second is the
sum of the point matrices at zero and at the $i$th unit vector.  The
third is the sum of the point matrices at zero, at the $i$th and $j$th
unit vectors, and at their sum.  This proves the first assertion.

To isolate a selector $s$ among $u$ distinct selectors, for each other
selector choose a coordinate on which it differs from $s$, and take
the product of the corresponding affine bit functions that equal one
at $s$ and zero at the chosen other selector.  This polynomial has
degree at most $u-1$ and evaluates to the indicator of $s$ on the
specified set.  Applying these interpolating linear functionals to a
relation among the evaluation vectors proves independence.
\end{proof}

\begin{lemma}[Label support]\label{lem:label-support}
Every $w\in W$ of rank at most $R_*$ has a representation
\begin{equation}\label{eq:label-decomposition}
 w=\sum_{s\in L}(p_s\otimes I)Z_s(p_s\otimes I)\trans,
 \qquad |L|\le R_*,
\end{equation}
where each $Z_s$ is symmetric, $(Z_s)_{ii}=(Z_s)_{0i}$, and
$\sum_{s\in L}\rank Z_s\le R_*$.
\end{lemma}
\begin{proof}
We first construct the label summands at a selector degree where
the pairing rank has stabilized.  We then show that any discrepancy
at a higher degree would contradict the rank bounds.

Choose a linear combination of point matrices representing $w$.
It defines a linear functional $\Lambda$ on the Boolean functions
of selector degree at most $2j_*$ and base degree at most two.
The matrix $w$ is the matrix of the pairing
$B(f,g)=\Lambda(fg)$ on functions that are base-linear and have
selector degree at most $j_*$.  Products and degrees here are in
the Boolean algebra, so each variable satisfies $x^2=x$.

Let $V_j$ denote the subspace with selector degree at most $j$,
and let $r_j$ be the rank of $B|_{V_j\times V_j}$.  The sequence
$r_j$ is nondecreasing and bounded by $R_*$.  Because
$j_*\ge 6R_*+4$, there is $j\ge4R_*+2$ such that
\begin{equation}\label{eq:flat-levels}
 r_j=r_{j+1}=r_{j+2}.
\end{equation}
Indeed, if no such three successive values occurred between
levels $4R_*+2$ and $j_*$, every two successive increments in
that interval would contain a strict increase, giving more than
$R_*$ increases.

Let $Q$ be the quotient of $V_{j+2}$ by the radical of its
restricted pairing.  The induced pairing on $Q$ is nondegenerate.
Equality of the first and last ranks in \eqref{eq:flat-levels}
means that the image of $V_j$ is all of $Q$, and that a vector of
$V_j$ pairing to zero with $V_j$ pairs to zero with $V_{j+2}$.
For a selector coordinate $s_a$, define $M_a[f]=[s_a f]$ using
$f\in V_j$.  This is well defined: if $[f]=0$, then for every
$g\in V_j$,
\[
 B(s_a f,g)=B(f,s_a g)=0,
\]
and $V_j$ spans $Q$.  The same identity proves that $M_a$ is
self-adjoint.  Replacing $s_bf$ by a representative in $V_j$ and
transferring $s_a$ to the other argument gives
\[
 B(M_aM_b[f],[g])=\Lambda(s_as_bfg).
\]
All representatives needed in this comparison lie in $V_{j+2}$.
Thus $M_aM_b=M_bM_a$ and, since $s_a^2=s_a$, $M_a^2=M_a$.

Commuting idempotents over $\F$ have a simultaneous eigenspace
splitting
\[
 Q=\bigoplus_{s\in L}Q_s,
 \qquad M_a|_{Q_s}=s_a I.
\]
For completeness, a single idempotent splits its space into kernel
and image; all the other commuting maps preserve both summands,
so iteration proves the assertion.  Distinct joint eigenspaces
are orthogonal, by self-adjointness of a coordinate on which their
labels differ.  Each nonzero summand has positive dimension, so
$|L|\le\dim Q\le R_*$.  Define $Z_s$ by pairing the projections
onto $Q_s$ of the affine base coordinates.  Then $Z_s$ is symmetric
and $\rank Z_s\le\dim Q_s$.

Repeated application of the $M_a$ gives the class of a base
coordinate times any selector monomial of degree at most $j$.
It follows that \eqref{eq:label-decomposition} reproduces all
moments of selector degree at most $2j$.  To verify the diagonal
identity separately on a label, use an interpolating selector
polynomial $e_s$ of degree at most $|L|-1$.  Its operator on $Q$
is the projection onto $Q_s$.  The Boolean identities give
\[
 B(z_i e_s,z_i e_s)
 =\Lambda(z_i e_s^2)
 =B(e_s,z_i e_s),
\]
which is $(Z_s)_{ii}=(Z_s)_{0i}$.  The interpolation degrees are
within $j$.  We have therefore obtained the required label blocks
and all moments through selector degree $2j$.

It remains to recover the higher-degree moments.
Extend the right side of \eqref{eq:label-decomposition} to the
full row and column indexing.  Its rank is at most
$\sum_s\dim Q_s\le R_*$.  By \Cref{lem:base-moments} it belongs
to $W$.  If its difference from $w$ were nonzero, choose a
nonzero difference moment of minimal selector degree $d'>2j$.
Write its selector monomial as the product on a set $U$ of size
$d'$, and keep the two affine base coordinates in that moment
fixed.  Index rows by the subsets $A\subset U$ of size
$\lfloor d'/2\rfloor$ and columns by the complementary monomials
$U\setminus A'$ with the same indexing.  The union
$A\cup(U\setminus A')$ is $U$ exactly when $A=A'$; otherwise
its degree is smaller than $d'$.  The difference on this square
submatrix is therefore the identity.  Both its row and column
degrees are at most $j_*$, whereas its rank is
$\binom{d'}{\lfloor d'/2\rfloor}>2R_*$.  This contradicts the
rank bound $2R_*$ for the difference and completes the proof.
\end{proof}

\subsection{Labels protected from the pins}

The decomposition separates the label blocks of a low-rank profile.
Projected pin directions are arbitrary vectors of $\cB$, so we need
a separate bound on the labels in their sparse expansions.
For an endpoint $i$ of a leaf, let $S_{i,e}^{\pm}\subset\cB$
be the projection of $D_{O,e}^{\pm}$ onto its individual primal
coordinate block.  These are projections, not intersections;
a stored pin may mix endpoints and channel coordinates.
The label support of a sparse expansion
$\sum_s p_s\otimes z_s$ is the set of labels with $z_s\ne0$.

\begin{lemma}\label{lem:sparse-pin-labels}
Put $L_0^{\rm lab}=R_*+14$.  At an endpoint, the union of the
nonzero label supports of all vectors in all $S_{i,e}^{\pm}$
that can be written as
$\sum_{s\in L}p_s\otimes z_s$ with $|L|\le L_0^{\rm lab}$
has size at most
\begin{equation}\label{eq:label-budget}
 B_0=2|\cE|K(R_*+14).
\end{equation}
\end{lemma}
\begin{proof}
For one projected space, choose a basis from its members admitting
such sparse representations, for the span of all these members.
There are at most $K$ chosen members.  Comparing any further sparse
member with its expression in this basis uses at most
$(K+1)(R_*+14)$ labels.  Their evaluation vectors are independent
by \Cref{lem:base-moments} and \eqref{eq:moment-parameters}.
Consequently every label of the further member occurs in one of
the chosen basis representations.  There are at most
$K(R_*+14)$ such labels per projected space.  Summing over signs
and components proves the stated, deliberately generous, bound.
\end{proof}

We will choose the witness points at this endpoint outside these
labels.  An expansion on at most $R_*+14$ labels of a vector in a
projected pin space can then have no nonzero summand at a chosen label.
For later finite exclusions set
\[
 B_*=10(B_0+10000g+1)
\]
and choose $b$ so large that
\begin{equation}\label{eq:selector-size}
 2^{b-2j_*}>100g(B_*+1).
\end{equation}
The part of the cut space supported entirely in the projected pins
will require separate scalar tests.  Define the effective test space
at endpoint $i$ by
\begin{equation}\label{eq:effective-space}
 C_i=\{x\in\cX:x_e\in S_{i,e}^+\otimes S_{i,e}^-
                  \text{ for every }e\in\cE\}.
\end{equation}
Fix a basis of this space on each leaf.  Since the total stored
pin dimension is at most $K$, the sum over components of the
individual projection dimensions in either sign is at most $K$.
Therefore
\begin{equation}\label{eq:effective-bounds}
 \dim C_i\le K^2,\qquad
 \sum_e\rank x_e\le K\quad(x\in C_i).
\end{equation}
A basis therefore reduces equality of linear functionals on $C_i$
to boundedly many scalar conditions.
If pins are enlarged, retain the old effective spaces and their
chosen basis data as well; they are subspaces of the new ones.

We now specify the point profiles used to build witnesses.
\begin{definition}[Atoms and flavors]\label{def:atom}
For a tag $l$, selector $s$, and base point $z$ allowed at $l$,
let $\mathfrak a_l(s,z)\in\cX$ be the cut profile represented
by $w_l=v(s,z)v(s,z)\trans$ and $w_t=0$ for $t\ne l$.
Thus its component is $v(s,z)v(s,z)\trans$ at every component
incident with $l$, and zero elsewhere.
Its diagonal bit is $q=\eta(w_l)$ and its role bit is
$a(\mathfrak a_l(s,z))$.
A generic flavor retains all base coordinates allowed at $l$.
A shared-only flavor sets all $O$ blocks to zero.  A pure flavor
sets $S$ to zero and retains a specified subset of the allowed
$O$ blocks, setting the others to zero.  Every flavor retains
both $\#$ and $Z$ in their entirety.
At a fixed label, a parameter test samples all retained bits
independently and uniformly.  Its tester summary is the list of
all individual quadratic tester values.  In addition to this
summary we may record the basis values of
$T(\mathfrak a_l(s,z),\cdot)|_{C_i}$.
\end{definition}

\subsection{Choice of the mixers}

For one atom, the scalar tests include the values of
$T(\mathfrak a_l(s,z),\cdot)$ on a basis of $C_i$.
Every nonzero combination of these tests is evaluation against
a nonzero profile of total
component rank at most $K$.  For two atoms with distinct labels,
we will also prescribe ordered $L,R$ evaluations.  The two parts
of the next lemma give the rank bounds used to control these
one-point and two-point tests.

Fix, for example,
\begin{equation}\label{eq:mixer-parameters}
 A=10(K+1)p_*(g^2+5),\qquad s_0=4\lceil A\rceil,
 \qquad J>100(s_0+K^2+1).
\end{equation}
These constants precede $r_0,h$ and $M_0$.

\begin{lemma}[Uniform mixing forms]\label{lem:generic-mixers}
For all sufficiently large $n$ one can choose the matrices
$L_j^{ef},R_j^{ef}$ and $M_a$ in \Cref{sec:geometry} so that:
\begin{enumerate}[label=\textup{(\roman*)}]
\item\label{item:mixer-unary}
For every nonzero $x\in\cX$ of total component rank at most $K$,
every tag, label and flavor, and every fixing of the non-$Z$
base coordinates, the quadratic function
$T(x,\mathfrak a_l(s,z))$ of the $n$ remaining $Z$ bits has polar
rank at least $2(J-s_0)$.  It uses a space of linear forms on $Z$
of bounded dimension, independent of $n$.  This space can be
chosen before the non-$Z$ coordinates are fixed.
\item\label{item:mixer-binary}
For two distinct labels and any two flavors, every nonzero parity
of the indexed $L_j^{ef},R_j^{ef}$ evaluations between the two
point inputs, allowing each matrix in either order or both,
and optionally either or both orders of $E$, has bilinear-part
rank at least $n/5$.  Here a nonzero parity involving an $L$ or
$R$ evaluation means that at least one such indexed ordered
evaluation occurs.  The same rank bound holds for the three
parities consisting of $E$ alone, its reverse alone, or their sum.
\end{enumerate}
The choices are made before any unit law is selected.
\end{lemma}
\begin{proof}
A symmetric rank-$r$ matrix on a $d_B$-dimensional coordinate
space can be written $ZHZ\trans$, where $Z$ has $r$ independent
columns and $H$ is a nonsingular symmetric $r$-by-$r$ matrix.
To verify the factorization, take a left inverse $L$ of the
chosen column basis $Z$ and put $H=LxL\trans$.  The projection
$ZL$ fixes $x$ on the left and, by symmetry, on the right;
hence $ZHZ\trans=x$.  Its rank forces $H$ to be nonsingular.
Neither $H$ nor the original matrix is required to have nonzero
diagonal.  For a fixed component rank tuple of total at most
$K$, counting $Z,H$ gives at most $2^{d_BK+K^2+K}$ choices.
There are at most $(K+1)^{|\cE|}$ such rank tuples.  Therefore
the number of profiles under consideration satisfies
\begin{equation}\label{eq:low-rank-profile-count}
 \#\{x\in\cX:\textstyle\sum_e\rank x_e\le K\}\le2^{An}
\end{equation}
for all sufficiently large $n$, since
$d_B=p_*(1+(g^2+3)n)$.

Choose all $L,R$ matrices independently and uniformly.  Fix a
nonzero profile and write $x_e=Z_eH_eZ_e\trans$, with
$r_e=\rank x_e$.  For a tag-$l$ atom the contribution from a
single indexed matrix uses the forward forms $Z_e\trans L v$
when $l\in f$, and the reverse forms $v\trans L Z_f$ when
$l\in e$, and the corresponding forms for $R$.
Regard these initially as separate formal variables.  Across
all indices their number is at most
\[
 m_0=4J(g-1)K.
\]
Each forward or reverse pair contributes a bilinear product
with nonsingular middle matrix $H_e$ or $H_f$.  These products
use disjoint formal variables.  Since $x\ne0$, for each $j$
there is at least one such nonzero product, so the formal
quadratic has polar rank at least $2J$.

We justify the distribution of their linear parts on $Z$ even
when forward and reverse restrictions overlap.  For a fixed
matrix $L$, specifying $Z_e\trans L$ and $LZ_f$ imposes at
most $r_er_f$ compatibility bits, namely their common
$Z_e\trans LZ_f$ block.  After restriction to the free $Z$
directions, the resulting list remains uniform on a linear
subspace of codimension at most $r_er_f$ in the space of all
lists.  Equivalently its law is dominated by $2^{r_er_f}$ times
the completely uniform law.  Over all matrices the domination
factor is at most $2^{2JK^2}$; only matrices with both relevant
incidences contribute a compatibility term.

For a uniform $m$-by-$n$ matrix, the probability of corank at
least $s_0+1$ is at most
$2^{m(s_0+1)-(s_0+1)n}$: choose $s_0+1$ independent row
relations and require all their $n$ evaluations to vanish.
If $m<s_0+1$ the probability is zero.  The preceding domination
therefore bounds the bad probability for a fixed profile, tag
and label by
\[
 2^{2JK^2+(s_0+1)m_0}\,2^{-(s_0+1)n}.
\]
There is no union over non-$Z$ values or flavors in this estimate:
they affect affine shifts, not these linear parts on the always
retained $Z$ block.  A union over at most $2^{An}g2^b$ choices
tends to zero.  On the complementary event the image of the
$Z$ inputs in the formal variable space has codimension at most
$s_0$.  Restricting a bilinear form to a subspace of codimension
$s_0$ reduces rank by at most $2s_0$: in a basis extending
that subspace one deletes $s_0$ rows and $s_0$ columns, each
deletion decreasing rank by at most one.  This proves
\ref{item:mixer-unary}.  The formal list itself supplies the
claimed bounded space of linear forms.  The term $T^0$ has no
$Z$ variables and does not affect this argument.

For \ref{item:mixer-binary}, distinct labels have independent
selector evaluation vectors.  Their two $Z$ direction spaces
are therefore disjoint.  Restricting a random matrix to these
two spaces in the two orders reads independent off-diagonal
blocks.  Any nonempty parity involving ordered $L,R$ evaluations
is consequently a uniform $n$-by-$n$ matrix on these directions,
even after all other matrices and $E$ have been fixed.

For the parities using only $E$, choose its matrices $M_a$
independently and uniformly.  On the two $\#$ direction spaces,
$E$ restricts to
\[
 A_{s,s'}=\sum_{a=1}^b(s_a+s'_a)M_a,
\]
which is uniform because $s\ne s'$.  The reverse gives
$A_{s,s'}\trans$, and their sum is alternating.  An
$\lfloor n/2\rfloor$ square block with disjoint row and column
indices is uniform also in this last case.  A uniform
$d_1$-by-$d_2$ matrix has probability at most
$2^{(d_1+d_2)r-d_1d_2}$ of rank at most $r$, by representing
such a matrix as a product of a $d_1$-by-$r$ and an
$r$-by-$d_2$ matrix.  Taking $r< n/5$ gives probability
$2^{-\Omega(n^2)}$ in all the preceding cases, including the
square block of side $\lfloor n/2\rfloor$.

There are at most $2^{2b}$ ordered label pairs and
$2^{4J|\cE|^2+2}$ parities to check.  These counts are fixed
in $n$.  The $Z$ and $\#$ restrictions also show that the
tester part of $E$ and the choice of flavor cannot reduce the
ranks just proved.  A union bound establishes both properties
simultaneously.  Fix one successful choice for each sufficiently
large $n$.
\end{proof}

\begin{remark}\label{rem:Z-reparametrization}
The raw frame law has an exact symmetry under a common
$\GL_n(2)$ change of the $Z$ input coordinates at one vertex,
using that change at every component and in both primal modes.
Indeed $E$ has no $Z$ entries, so the prescribed self-Grams and
all frame injectivities are preserved; the transformation is a
bijection of the finite uniform frame space.  Tester summaries
are unchanged.  This is a symmetry of the raw law, not a claim
that $T^1$ is invariant.  Under a change of variables in a
parameter test it moves the bounded row-bit space in
\Cref{lem:generic-mixers}, part~\ref{item:mixer-unary} while preserving
the unfiltered empirical mass of any test depending only on the
resulting key and on unchanged tester data.
\end{remark}

Choose $r_0$ only after the preceding parameters and the bounded
rank requirements in \Cref{sec:realization,sec:phases}; retain
$h=1000r_0$.  Finally choose $M_0$ large enough for every stated
nominal-coefficient and scalar-record cost to be the required
small fraction of $N=M_0n$.  None of the finite-test accuracies
used later changes these choices.
